\section{模型的检验与灵敏度分析}
\label{sec:6}

\subsection{模型检验}
\label{sec:6.1}

全文检验汇总于表~\ref{tab:checks}。「互相验证」用于同一问题、不同假设的方法对，判据取排序或数值的一致性；上下游分工衔接的方法之间不构成验证关系。

\begin{table}[H]
\centering\footnotesize
\caption{全文检验清单}
\label{tab:checks}
\begin{tabularx}{\textwidth}{@{}c p{3.85cm}p{1.55cm}p{2.25cm}X@{}}
\toprule
问 & 检验对 & 类型 & 判据 & 结果 \\
\midrule
一 & Huber 主分 vs 均值/PCA/TOPSIS & 方法敏感性 & 样本 Spearman & $0.9936/0.8166/0.9588$ \\
一 & A1 vs A2/A3 & 扩展检验 & KS 的 $p$ 值 & $0.313/0.220$，未检出分布差异 \\
一 & 六来源候选跨集 & 规则复检 & 候选比例与方向 & $5.00\%/4.73\%/4.99\%$，主导方向有变化 \\
一 & clr+Huber vs 裸份额 OLS vs GBDT & 1M 留出 & 逐域/目标 $R^2$ & $0.772/0.260$、$0.654/0.418$、$0.923/0.870$ \\
一 & A1$\to$A18 质量桥接 & 五折＋六映射域 & $R^2$、MAE & $0.4106$、$0.0365$（$n=6$） \\
一 & 1M $\to$ 60M $\to$ 1B & 跨尺度（1B 独立） & 目标 Spearman & $0.523/0.467/0.677$ \\
一 & 尺度修正（幂律主口径） & 1B 独立检验 & 逐域 MAE & $2.953\to0.770$，降幅约 $74\%$ \\
一 & A12–A15 估算表 & 一致性对照 & 排序 & 与实测相反，作反例 \\
二 & 有地板 vs 无地板 & 嵌套 $F$ 检验 & $F$、$p$ & $F=307.11$，$p<10^{-15}$ \\
二 & $\kappa_N=\kappa_D$、$\kappa_D=0$ & 嵌套 $F$ 检验 & $F$、$p$ & $57.19$（$3.4\times10^{-13}$）拒绝；$0.63$（$0.43$）不拒绝 \\
二 & B7 新增 90 点 & 留出（插值） & RMSE、偏差 & $0.0531$、$+0.0219$；噪声代理 $0.0458$ \\
二 & 解析式 vs 数值解 & 互相验证 & 相对误差 & $<10^{-14}$（替代倍数）、$<10^{-8}$ \\
二 & $Q=1$ 切片 vs $L_0$ & 退化核对 & 残差区间 & $+0.0036$，$[-0.0120,0.0193]$ \\
二 & B4/B5/B2/B10 & 外部趋势对照 & Spearman & $0.983/0.959/0.788/1.000$，水平不可比 \\
二 & B11/B12 来源审计 & 元数据核对 & 覆盖、$D$ 记账 & 步数覆盖 $100\%$，中位比 $0.9986$ \\
三 & 数值解 vs 闭式 $N^*$ & 互相验证 & 相对误差 & $1.54\times10^{-7}$ \\
三 & 求解器 vs 粗网格 & 互相验证 & Loss 差 & 最大差 $1.80\times10^{-5}$ \\
三 & 剖面导数 & 有限差分 & 相对误差 & $4.1\times10^{-10}$ \\
四 & C8 重算 vs C1 榜单 & 口径核验 & Pearson、MAE & $0.9893$–$0.9995$；MATH $0.8009/3.4974$ \\
四 & S 形 vs 线性 vs 常数 & 留一 & RMSE & $8.079/8.283/9.514$ \\
四 & 分解三口径 & 方法敏感性 & 技术占比极差 & $32.81$ pp \\
四 & 预测回测 & 事后回测 & 2025.25 误差 & M1b $-0.784$；朴素线性 $+2.651$ \\
\bottomrule
\end{tabularx}
\end{table}

\textbf{问题一。} 检验分三层：A2/A3 为冻结评分的扩展检验，A6/A7 为配比模型的同尺度留出，A8–A11 为跨尺度检验，其中 A10/A11（1B）完全不参与估计、A8/A9 承担尺度修正标定。GBDT 依据 A6/A7 表现选为交叉评估器，该留出不再重复充当候选收益的独立验证。A12–A15 为估算/外推表，其隐含的大尺度排序反转与实测排序保持相矛盾，作反例使用。

\textbf{问题二。} 质量参数只用 B6（半合成）估计，留出只用 B7 新增 90 点（$Q$ 方向插值）。B2 为半合成、B3 由 B1 插值派生、B10 为公式估算，二者用于核对实现与趋势；B4 的 57 点剔除 13 个近邻后 44 点结果稳定（Spearman $0.975$）。B8 与 B6/B7 方向相反（$64.26\%$ 低于 $E$），作反例。B11/B12 审计确认 B1 的 8 个名义规模、每模型 147 步均被覆盖，B1 训练设定真实、损失按公开公式生成。

\textbf{问题三。} 数值最优解与闭式 $N^*$、求解器与独立粗网格、剖面解析导数与有限差分三组互验，相对误差分别为 $1.54\times10^{-7}$、最大 Loss 差 $1.80\times10^{-5}$、$4.1\times10^{-10}$，预算残差与份额和误差均小于 $10^{-9}$；锚点映射与单域越界均按规范标记。

\textbf{问题四。} C8 重算与 C1 为同批数据的口径核验（MATH 差异待解释），桥接为同域留一，分解差异为方法敏感性，回测为事后回测且重估样本未严格截在起点前。问题四目前没有独立的时间外检验层，结论归属拟合层与条件外推层，这一边界在 \ref{sec:7.2} 说明。

\subsection{灵敏度与稳健性分析}
\label{sec:6.2}

\textbf{问题一。} 六来源阈值取 P90/P95/P97.5 时候选比例约 $10.00\%/5.00\%/2.50\%$（A2 $9.19\%/4.73\%/2.00\%$，A3 $9.74\%/4.99\%/2.44\%$），阈值直接决定人工复核工作量。折扣对照分 $Q_{\rm v2}$ 与主分秩相关达 $0.9547/0.9207/0.9395$、域级排序不变。单指标极端扰动下主分变动约小 $30\%$；13 域等权与 pile\_cc 单域两目标给出不同候选，评价目标在搜索前固定；质量交互项仅 6/13 域显著。

\textbf{问题二。} 去地板后质量 $+0.1$ 在千亿参数处的等效倍数由约 $3.0$ 降到 $1.3$–$1.5$，方向不变；$\kappa_N$、$k_0$ 的 $95\%$ 区间为 $[0.1264,0.1568]$、$[0.2013,0.2554]$，$\kappa_D$ 区间含 0；主式形状代价为留出 RMSE $0.0531$（线性对照 $0.0416$、噪声 $0.0458$）。$C=10^{19}$ 时提质是否划算取决于成本函数，$C\ge10^{20}$ 后三种 $g(Q)$ 下提质都更有效。

\textbf{问题三。} 敏感性覆盖三类成本、C7 五档上下文、$\eta$ 三档、P2 重抽样、两种质量映射与配比不确定性，以及 \texttt{m=1}、H$_{\rm tot}$、无地板对照。与预算和规模无关的解析结论是 $L_{ctx}^{\rm crit}=30000$ token；仅含 P2 时，离开下界的 bootstrap 区间为指数型 $[19.038,19.084]$、幂函数型 $[19.503,19.560]$、对数型 $[18.539,18.620]$。Q3-9 已将 P1 域质量与配比 bootstrap 与 P2 联合传播，离开下界区间外扩约 $0.003$ dex，三种成本函数的阈值排序不变。

\textbf{问题四。} Medium 层权重在 $0.25$–$1.0$ 时 $S(1.85)$ 变化不足 $0.4$ 分，仅用 High 层时上平台降到 $6.349$，上升段依赖 Medium 层。C3 权重取 $0.2$、$1.0$ 时技术率为 $0.499$、$0.205$ dex/年（不含 C3 为 $0.781$）。换窗口后 M1b 技术占比为 $7.06\%$、$8.18\%$，「规模为主」稳定。预测对后训练增益是否延续最敏感（约 $6.8$ 分），大于算力情景差异（约 $2$ 分）。
